\ExamPart{Trắc nghiệm đúng sai}{2}{Thí sinh trả lời từ câu 1 đến câu 2. Trong mỗi ý a), b), c), d) ở mỗi câu học sinh chọn đúng hoặc sai.}

\NQuestion Xét tam thức bậc hai $f(x)=x^2-6x+8$ và các khẳng định sau:
\begin{SubQuestions}
  \TFItem{true}{Phương trình $f(x)=0$ có hai nghiệm phân biệt là $2$ và $4$.}
  \TFItem{true}{Trục đối xứng của parabol $y=f(x)$ là đường thẳng $x=3$.}
  \TFItem{false}{Bất phương trình $f(x)<0$ có tập nghiệm là $(-\infty;2)\cup(4;+\infty)$.}
  \TFItem{true}{Giá trị nhỏ nhất của $f(x)$ bằng $-1$.}
\end{SubQuestions}
\SolutionBlock{Ta có
\[
f(x)=x^2-6x+8=(x-2)(x-4).
\]
Vậy phương trình có hai nghiệm $2,4$, trục đối xứng là
\[
x=\dfrac{2+4}{2}=3.
\]
Vì $a>0$ nên $f(x)<0$ khi $2<x<4$, do đó mệnh đề c) sai.
Mặt khác
\[
f(3)=9-18+8=-1,
\]
nên giá trị nhỏ nhất bằng $-1$.}

\NQuestion Trong mặt phẳng tọa độ $Oxy$, cho $A(1;2)$, $B(5;4)$ và đường thẳng $d:x-y+1=0$. Xét các khẳng định sau:
\begin{SubQuestions}
  \TFItem{true}{Tọa độ vectơ $\overrightarrow{AB}$ là $(4;2)$.}
  \TFItem{false}{Trung điểm của đoạn thẳng $AB$ là $M(2;3)$.}
  \TFItem{true}{Đường thẳng $d$ có một vectơ pháp tuyến là $(1;-1)$.}
  \TFItem{false}{Điểm $A$ thuộc đường thẳng $d$.}
\end{SubQuestions}
\SolutionBlock{Ta có
\[
\overrightarrow{AB}=(5-1;4-2)=(4;2),
\]
nên a) đúng.
Trung điểm của $AB$ là
\[
M\left(\dfrac{1+5}{2};\dfrac{2+4}{2}\right)=(3;3),
\]
nên b) sai.
Đường thẳng $d:x-y+1=0$ có vectơ pháp tuyến $(1;-1)$ nên c) đúng.
Thay tọa độ điểm $A(1;2)$ vào $d$ được $1-2+1=0$, nên thực ra $A\in d$.
Vì thế mệnh đề d) sai.}
